\section*{System Requirements}
\addcontentsline{toc}{section}{System Requirements}

\subsection*{Functional requirements}

The functional requirements come from sections 2, 2.1 and 3 of the project
guideline. The two marked \emph{ours} are additions the guideline allows in
section 3. Every later artefact (use cases, scenarios, tests) refers to these
IDs, so they are never renumbered.

\begin{longtable}{@{}p{1.3cm}>{\raggedright\arraybackslash}p{8.6cm}p{2.3cm}p{1.8cm}@{}}
\toprule
\textbf{ID} & \textbf{Requirement} & \textbf{Guideline} & \textbf{Source}\\
\midrule
\endhead
FR-1 & Search for flights by origin, destination, date and airline. & 3.1 & guideline\\
FR-2 & Display available seats, flight times and prices. & 3.1 & guideline\\
FR-3 & Create new reservations. & 2, 3.2 & guideline\\
FR-4 & Cancel or modify existing reservations. Customers cancel, flight agents modify. & 2, 2.1, 3.2 & guideline\\
FR-5 & Generate a reservation confirmation, shown on screen and saved as a \texttt{.txt} file. & 3.2 & guideline\\
FR-6 & Create customer profiles, and view and edit customer details. & 2, 3.3 & guideline\\
FR-7 & Add flights, update flight details and remove outdated flights. Administrator only. & 2, 3.4 & guideline\\
FR-8 & Manage routes and aircraft, and update schedules. & 2.1 & guideline\\
FR-9 & Record simulated payment information and keep payment records. & 2, 3.5 & guideline\\
FR-10 & View reservation history and flight schedules. & 2, 2.1 & guideline\\
FR-11 & Store flight and reservation data in a local SQLite database. & 3.6 & guideline\\
FR-12 & Show a monthly promotion, picked at random the first time the application opens in a month and shown once that month. & 3.7 & guideline\\
FR-13 & Log in with one of three roles (customer, flight agent, administrator); the available features depend on the role. & 2.1 & ours\\
FR-14 & Select a seat when making a reservation. & 3.1 & ours\\
\bottomrule
\end{longtable}

\subsection*{Non-functional requirements}

Section 4 of the guideline names usability, portability, maintainability and
scalability. The assessment criteria add reusability and reliability, and we
add security. Each quality has a criterion, the part of the design
that meets it, and the automated tests that check it.

\begingroup
\raggedright
\begin{longtable}{@{}>{\raggedright\arraybackslash}p{2.8cm}>{\raggedright\arraybackslash}p{4.1cm}>{\raggedright\arraybackslash}p{3.5cm}>{\raggedright\arraybackslash}p{4.0cm}@{}}
\toprule
\textbf{Quality} & \textbf{Measurable criterion} & \textbf{How the design meets it} & \textbf{Checked by}\\
\midrule
\endhead
Usability & Every screen of a role is a tab or a button of that role's start screen. A bad input shows one plain sentence that names the field and the problem, never a stack trace. Enter submits the default button and Tab uses Swing's own layout order: no screen overrides it. & A top bar with one tab per task of the role. Shared helpers parse input and phrase errors, so every form behaves the same. & \footnotesize \texttt{NavModelTest}\newline \texttt{KeyboardTraversalTest}\newline \texttt{ScreenDefaultsTest}\newline \texttt{InputParserTest}\newline \texttt{ErrorMessagesTest}\\
Portability & Running the jar with the argument \texttt{-{}-selftest} prints \texttt{SELFTEST OK} and exits with 0. The jar needs only Java 17. The database is created on first run, with no installer. & Standalone Swing application, embedded SQLite, one jar that holds every dependency. & \footnotesize \texttt{SelfTestTest}\newline \texttt{MainFrameSmokeTest}\newline the build workflow runs on Linux with Java 17; window tests need a display and are skipped there\\
Maintainability & 0 Checkstyle violations on every build. The layers depend only downward, and ArchUnit tests fail the build if they do not. & Presentation, business (control and entity) and data packages. Class names match the class diagram. & \footnotesize \texttt{ArchitectureTest}\newline Checkstyle in every \texttt{mvn verify}\\
Scalability & Replacing SQLite changes only the data package: presentation and business classes have no \texttt{java.sql} or \texttt{Sqlite*} dependency. Controllers contain no Swing code. & DAO interfaces with package-private SQLite classes behind a factory. Controllers take their DAOs through the constructor. & \footnotesize \texttt{ArchitectureTest}\newline \texttt{DaoFactoryTest}\\
Reusability & Every entity and controller runs without Swing and without a display. The controller tests and the end-to-end test pass on a machine with no screen. Entities need no database. & Entities know nothing about the other layers. Controllers receive their DAOs and a clock instead of creating them. & \footnotesize \texttt{BookingControllerTest}\newline \texttt{PaymentControllerTest}\newline \texttt{EndToEndFlowTest}\newline \texttt{AppContextTest}\\
Reliability & A seat cannot be reserved twice: a second active reservation for the same flight and seat is rejected by the database, and a test proves it. A declined payment never leaves a confirmed reservation. & A partial unique index is the last line of defence. The reservation DAO and the \texttt{BookingController} turn it into a clear message. & \footnotesize \texttt{SqliteReservationDaoTest}\newline \texttt{BookingControllerTest}\newline \texttt{PaymentControllerTest}\\
Security & Passwords are stored only as PBKDF2 hashes with a random salt. Only the last 4 card digits are stored, and a card number is never logged. & A password hasher in the utility package. The payment details class keeps only the masked card. & \footnotesize \texttt{PasswordHasherTest}\newline \texttt{PaymentDetailsTest}\newline \texttt{SqlitePaymentDaoTest}\\
\bottomrule
\end{longtable}
\endgroup
